\section{Introduction}

\IEEEPARstart{T}{he} electrification of heating in cold-climate regions concentrates
distribution-system stress into a few winter hours~\cite{Knittel2025}. In the
province of Qu\'ebec, where electric space heating is the dominant heating
technology, the system peak is driven almost entirely by simultaneous resistive
and heat-pump heating on the coldest mornings and evenings, and field trials
confirm that aggregated electric heating sharply raises coincident peak
demand~\cite{Love2017}. As distributed energy resources (DERs) proliferate and
heating loads grow, these synchronized peaks threaten thermal limits on feeder
trunk conductors and depress voltages far from the substation, eroding the
hosting capacity of low-voltage networks~\cite{Protopapadaki2017,Edmunds2021}.
Demand response that shifts flexible loads to relieve these peaks is the
established mitigation~\cite{Siano2014}. Because the flexibility lives behind
thousands of meters and not at a single controllable node, centralized dispatch
scales poorly; decentralized coordination, in which each home solves a small
local problem and exchanges only an aggregate signal, is the natural fit.

The alternating direction method of multipliers (ADMM) \cite{boyd2011admm} is a
workhorse for this setting: its sharing form lets a coordinator flatten a common
resource (here, total feeder load) while each agent retains private control of
its own heating schedule. ADMM-based demand response is well
studied~\cite{Tsai2017,Nguyen2018}, yet two gaps undermine its real-world
credibility.

The first gap is \emph{fragility to communication failures}. Multi-agent
coordination implicitly assumes every agent reports on every iteration; when a
fraction of agents fall silent, the aggregate the coordinator optimizes is no
longer the aggregate that materializes on the wire. How badly the coordinated
outcome degrades, and how an operator might recover it by forecasting the
missing contributions, is rarely quantified end to end~\cite{Brown2017,Xu2019}.

The second gap is \emph{physical validity}. The literature typically stops at
the coordinated aggregate kilowatt signal and reports peak or peak-to-average
reductions. It almost never asks the operational question that matters: is the
coordinated aggregate actually feasible on the feeder? A flatter aggregate is
worthless if it still overloads a trunk line or sags a remote bus below limits.

This paper closes the loop between coordination and feeder physics. We adapt the
sharing-ADMM-plus-imputation methodology from a prior multi-agent thesis and
\emph{extend} it with an AC power-flow validation stage, turning a purely
signal-level study into a network-validated one. Our contributions are:

\begin{itemize}
  \item A comfort-aware sharing-ADMM coordinator for deferrable cold-climate
  electric heating that flattens the feeder aggregate subject to a per-home
  indoor-temperature comfort penalty and a daily-energy conservation constraint,
  so load is shifted rather than curtailed.
  \item A machine-learning imputation stage that reconstructs non-responsive agents
  so coordination degrades gracefully as the communication-failure fraction grows;
  comparing three ML families (regularized-linear, tree-ensemble, and
  instance-based) against naive and no-imputation baselines, we find that
  \emph{whether} one imputes matters far more than which method.
  \item An AC power-flow validation loop that closes the coordinated aggregate back
  onto feeder physics, quantifying the restored transformer and voltage feasibility
  and its breakdown under a swept non-responsive fraction with Monte-Carlo risk
  bands.
\end{itemize}

The remainder of the paper is organized as follows.
Section~\ref{sec:methodology} is preceded by the related work and the system
model; Section~II surveys ADMM-based demand response, multi-agent resilience,
telemetry imputation, and distribution power flow, and states our delta against
each. Section~III formalizes the agents, the feeder, and the
communication-failure model. Section~\ref{sec:methodology} develops the
sharing-ADMM updates, the ML imputer, and the power-flow validation loop.
Section~V details the case-study configuration, Section~VI reports the aggregate,
robustness, and network-feasibility results, and Section~VII concludes.
